\textbf{Construct validity.}
Коэффициенты $k$, $a$ и $h$ агрегируют качество постановки, возможности
инструмента, проверку и исправление результата. Один и тот же численный профиль
может скрывать разные механизмы, а makespan не измеряет качество кода,
понимание системы, стоимость вычислений и последующие дефекты. Термины
«синхронный» и «концентрация» в ранних экспериментальных гипотезах оказались
недостаточно операциональными; соответствующие выводы не включены в число
подтверждённых.

\textbf{Internal validity.}
Точные серии используют детерминированные длительности, целочисленную сетку и
один решатель. Ошибки реализации снижаются независимыми проверками ограничений,
статусами оптимальности, замороженными входами и воспроизводимыми manifests,
но вычислительный эксперимент подтверждает свойства реализации формальной
модели, а не истинность её допущений о реальном процессе.

\textbf{Conclusion validity.}
Точные конечные сетки не доказывают универсальные закономерности за пределами
выбранных диапазонов. Числа 30/90, 12/13 и 96/100 описывают именно исследованные
матрицы, а не частоты в генеральной совокупности проектов. В больших сериях
EXP-07 и части EXP-08 оценивается фиксированная политика; её эффекты нельзя
переносить на $T^*$ без отдельного точного доказательства.

\textbf{External validity.}
Большинство DAG, длительностей и функций $C(P),\gamma(P)$ синтетические и не
калиброваны на продольных данных команд. Рассматривается один разработчик,
однородные агентные потоки и локальная блокировка с удержанием потока. Другие
планировщики могут переиспользовать ожидающую мощность, несколько людей могут
обслуживать запросы параллельно, а реальные задачи раскрывают зависимости и
повторную работу динамически.

\textbf{Parameter uncertainty.}
EXP-06 проверяет только заранее заданные случайные и неблагоприятные семейства
ошибок. Наблюдаемая устойчивость ранжирования не является гарантией против
произвольной совместной ошибки оценок. Перед практическим решением необходимы
локальная калибровка, интервальный анализ и повторное планирование по мере
появления фактических данных.
